% Extracted from attraction_wfigs.docx; see PLAN.md and tools/extract_pilot.py.
\section*{Chapter 2: The Fathers of Gravity}
\begin{description}[style=multiline,leftmargin=2em,labelwidth=1.5em]

% DOCX paragraph 2017
\item[1.]\hypertarget{note:02-01}{}In a reference system in which the two bodies have a
position given by the vectors\index{vector} \({\overrightarrow{r}}_{1}\) and
\({\overrightarrow{r}}_{2}\), their separation is given by the vector difference between
the two, \({\overrightarrow{r} = \overrightarrow{r}}_{1} - {\overrightarrow{r}}_{2}\),
whose magnitude (or modulus) is indicated as \(\left| \overrightarrow{r} \right| =\)
\emph{r}.

% DOCX paragraph 2018
\item[2.]\hypertarget{note:02-02}{}The second of arc is a unit of measurement of angles
used both in navigation and astronomy and represents an alternative way of expressing
the amplitude of the angle span by an arc. In practice, in addition to measuring the
amplitude of an angle in degrees, a set of units we are familiar with (but one that is
not convenient in navigation or astronomy), smaller units are introduced, namely, the
\emph{minutes} of arc and the \emph{seconds} of arc. Therefore, an arc of one degree is
divided into 60 minutes (each corresponding to 0.017 degrees) or 3,600 seconds (each
corresponding to 0.00028 degrees). Using degrees rather than arcseconds, we would say
that Mercury's perihelion\index{perihelion} changes by only 0.0121 degrees every
century!

\end{description}
